% \documentclass{report}

% \usepackage{graphicx}
% \usepackage{dcolumn}
% \usepackage{bm}
% \usepackage{amsfonts, amssymb, amsmath}
% \usepackage{braket}

%\newcommand{\ket}[1]{\textstyle \left| #1 \right\rangle \displaystyle}
%\newcommand{\bra}[1]{\textstyle \left\langle #1 \right| \displaystyle}
%\newcommand{\braket}[2]{\textstyle \left\langle #1 \left|  #2 \right\rangle\right. \displaystyle}
%\newcommand{\im}{\mathrm{i}}
%\newcommand{\e}{\mathrm{e}}
%\newcommand{\pwrt}[2]{\frac{\partial #1}{\partial #2}}
%\newcommand{\diff}{\mathrm{d}}

%\begin{document}

\subsection{Introduction of a variational basis}
One of the concerns with the variational method is selection of the
wavefunction. The wavefunction must reflect the form of the actual wavefunction,
though this form may not be known in any detail. In this work, we consider
replacing a single variational wavefunction with a "variational basis". In this
section, we will introduce, define, and refine this idea. 

To start with, we select a set of functions such that we may expand the envelope
function of a valley as
\begin{equation}
  \Phi^{(q)}_{\bm{r}}
  \equiv
  \sum_{k} c_{k} \Phi_{k\bm{r}}.
\end{equation}
We wish to consider the bound states of our system. As such, we can assume that
the function $F^{(z)}$ is part of the
$L^{2}_{\mathbb{C}}\left(\mathbb{R}^{3}\right)$ function space. It also worth
noting that this function space is also a separable Hilbert space. We seek a set
of orthonormal functions that span this separable Hilbert space. In this work,
we explore two possibilities for our complete basis. While a completely rigorous
mathematically sound interpretation of the following statement is beyond the
scope of this work, I'll state formally that the countable eigenstates of the
Hamiltonian of the isotropic harmonic oscillator form a complete basis for the
separable Hilbert space. As such, an obvious candidate for our orthonormal basis
emerges from the energy eigenstates of the isotropic 3-D harmonic oscillator,
namely
\begin{equation}
  \varphi_{klm;\eta;\bm{r}}
  \propto
  L_{k;\rho^{2}}^{\left(l+\frac{1}{2}\right)}
  \mathrm{e}^{-\frac{\rho^2}{2}}
  \rho^{l}
  Y_{lm;\chi\phi}.
\end{equation}
Above, $L_{k}^{\left(\alpha\right)}$ denotes the generalized Laguerre
polynomials with $k$ being a non-negative integer, and with $Y$ denoting the
spherical harmonics such that $l$ is a non-negative integer and $m$ is an
integer such that $\left|m\right| \le l$. Our polar spherical coordinate $\chi =
z/r = \cos\left(\theta\right)$ is an alternate polar spherical coordinate we
find more convenient to use in this work than $\theta$ directly. As well, $\rho
= r/\eta$ is the radial coordinate rescaled to be dimensionless with $\eta$ as a
unitful constant shared by each state in the basis. We pick this radial scaling
$\eta$ to generate the best representation of our wavefunction. In the isotropic
harmonic oscillator, $\eta = \sqrt{\frac{\hbar}{m\omega}}$ with mass $m$ and
angular frequency $\omega$. In our system, this constant $\eta$ may be optimized
to generate the lowest expectation energy for our system, thus giving it a
variational character.

More generally, we expect the wavefunction of an isolated donor to display some
manner of near spherical or ellipsoidal symmetry. As such one of the building
blocks of our basis set will be the aforementioned spherical harmonics, complete
on $L^{2}_{\mathbb{C}}\left(S^{2}\right)$ where $S^{2}$ is the unit sphere. A
further important observation as we design our basis involves the tensor product
of Hilbert spaces. To wit, for Hilbert spaces $\mathcal{H}_{1}$ and
$\mathcal{H}_{2}$ with orthonormal bases $\left\{\phi_{i}\right\}$ and
$\left\{\eta_{j}\right\}$ respectively, it is true that
$\left\{\phi_{i}\otimes\eta_{j}\right\}$ is an orthonormal basis for
$\mathcal{H}_{1}\otimes\mathcal{H}_{2}$. Articulating the rigorous mathematical
definition of the tensor product of Hilbert spaces is beyond the scope of this
document. It is sufficient for our purposes to keep in mind that we can
construct a basis for $L^{2}_{\mathbb{C}}\left(\mathbb{R}^{3}\right)$ by noting
that this function space is isomorphic to the tensor product space
$L^{2}_{\mathbb{C}}\left(S^{2}\right)
\otimes
L^{2}_{\mathbb{C}}\left(\mathbb{R}^{+},r^2\right)$.

% https://aalexan3.math.ncsu.edu/articles/tensor.pdf
% https://math.stackexchange.com/questions/1653561/how-to-construct-a-complete-set-in-l2-mathbbr3-starting-with-the-spheric

To find a complete basis on $L^{2}_{\mathbb{C}}\left(\mathbb{R}^{3}\right)$
then, we only need to find a complete basis on
$L^{2}_{\mathbb{C}}\left(\mathbb{R}^{+},r^2\right)$. We've already seen such a
set of basis functions emerge in the isotropic harmonic oscillator, namely  
\begin{equation}
  H_{kl;r}
  \propto
  L_{k;r^{2}}^{\left(l+\frac{1}{2}\right)}
  \mathrm{e}^{-\frac{r^2}{2}}
  r^{l}
\end{equation}
However, in addition to this basis, we also want to consider a basis with Slater
exponential decay. The formula for the wavefunctions of the isotropic harmonic
oscillator includes generalized Laguerre polynomials defined in the conventional
way\cite{LaguerrePolynomials}. For our purposes, I note that the Laguerre
polynomial $L_{k;x}^{\left(\alpha\right)}$ is complete over the Hilbert space
$L^{2}\left(\mathbb{R}^{+},x^\alpha\mathrm{e}^{-x}\right)$. For the harmonic
oscillator eigenstates with $x = r^{2}$ and $\alpha = l + \frac{1}{2}$, I note
that $x^{\alpha}\e^{-x}\,\diff{x} \propto r^{2l+2}\e^{-r^{2}}\,\diff{r}$. As
such, the above statement of completeness for the Laguerre polynomials is
equivalent to the completeness of $H_{kl;r}$ on
$L^{2}\left(\mathbb{R}^{+},r^{2}\right)$. 

To arrive at our Slater-decay basis we will once more introduce the associated
Laguerre polynomial as our building block. This time we shall take $x = r$.
We introduce then
\begin{equation}
  G_{kl;r}
  \propto
  L_{k;r}^{\left(2l+2\right)}
  \mathrm{e}^{-\frac{r}{2}}
  r^{l}
\end{equation}
As $L$ above is complete over
$L^{2}\left(\mathbb{R}^{+},r^\alpha\mathrm{e}^{-r}\right)$, then it can be
shown that $G$ is complete over $L^{2}\left(\mathbb{R}^{+},\rho^{2}\right)$. As
a final aside, we took $\alpha$ above to be such that $\alpha = 2l + 2$.
Further discussion on why this choice has been made will be discussed in
the \href{https://gitlab.com/cantankerous-impurity/Documentation/-/tree/core/SUNY--Buffalo/Dissertation/Supplementary%20Materials?ref_type=heads}{Supplementary Materials}.

In short, the two forms above for the radial part of our basis can be combined
and generalized for all positive integer $W$, though we are only concerns with
$W = 1,2$. Appropriately normalized, the combined form is
\begin{equation}
  F_{kl;\eta;r}^{\left(W\right)}
  =
  \sqrt{\frac{W}{\eta^{3}}}
  L_{k;\rho}^{\left(\alpha\right)}
  \mathrm{e}^{-\frac{\rho}{2}}
  \rho^{\frac{l}{W}},
  \quad\quad
  \rho = \left(\frac{r}{\eta}\right)^{W}
\label{naive_wf:radial_part}
.
\end{equation}
with $\alpha = \frac{2l+3-W}{W}$. With $W = 1,2$ the above function will be the
radial part of our basis. Above, we assume that the Laguerre polynomials are
defined conventionally and orthonormalized such that
\begin{equation}
\int\diff{\rho}\
L_{k_{0};\rho}^{\left(\alpha\right)}L_{k_{1};\rho}^{\left(\alpha\right)}
\rho^{\alpha}
\e^{-\rho}
  =
    \delta_{k_{0}k_{1}}
.
\end{equation}

Thus we may write the complete formula for our naive basis states with positive
real $\eta$ acting as a variational parameter.
\begin{equation}
  \Phi_{klm;\eta;\bm{r}}^{\left(W\right)}
  =
  F_{kl;\eta;r}^{\left(W\right)}
  Y_{lm;\chi\phi}.
  \label{basis-1}
\end{equation}
In the above formula, we assume appropriate normalization factors have been
applied to the radial function $F$ and the spherical harmonics $Y$ such that the
resulting basis state $\Phi$ is normalized.

We note that the envelope function with $W = 1$ yielding Slater-exponential
decay bares a strong resemblance to the bound states of the hydrogen atom. As we
consider using our basis states to represent the energy eigenstates of our
system, this is an important point to consider. Of note, there are two points of
divergence between the hydrogenic states and our Slater basis states.

First, we want to be sure that our basis can conveniently represent the
hydrogenic eigenstates. The hydrogenic wavefunctions have the form
\begin{equation}
  \Psi_{nlk}
  \propto
  L_{n-l-1;\rho}^{\left(2l+1\right)}
  \mathrm{e}^{-\frac{\rho}{2}}
  \rho^{l}
  Y_{lm;\chi\phi}
  \quad\quad
  \rho 
  =
  \frac{r}{\eta_{n}}.
\end{equation}
It is trivial to reindex the wave function such that $n = k + l + 1$ and show
that the two set of indices are isomorphic between the hydrogenic states and our
Slater basis states. We note that the Laguerre parameter differs between the
hydrogenic and Slater basis, being $\alpha_{\text{hydro}} = 2l+1$ and
$\alpha_{\text{Slater}} = 2l+2$ respectively. Despite this, it is the case
that $L_{k}^{\left(\alpha\right)} = L_{k}^{\left(\alpha+1\right)} -
L_{k-1}^{\left(\alpha+1\right)}$. In light of the above facts, it is clear that
each hydrogenic wavefunction can be represented by, at most, two states from our
basis given the correct variational scaling factor $\eta$.

Secondly, however, note that the scaling factor $\eta_{n}$ for each hydrogen
state depends on the quantum numbers of the states. It goes as
\begin{equation}
  \eta_{n} = \frac{n}{2}r_{0}
\end{equation}
where 
\begin{equation}
  r_{0} = \frac{4\pi\epsilon_{0}\hbar^{2}}{me^{2}}
\end{equation}
We see a conflict between our complete basis with a single parameter $\eta$ and
the bound hydrogenic states.  A larger orthonormal basis that tends towards
completeness can represent any desired state in principle, of course. In our
experience, however, the computational cost from expanding our orthonormal basis
does not see a worthwhile increase in that basis's power to represent more than
a single desired state.

To resolve this, we introduce the concepts that we will call in this document
the \emph{subbasis} and the \emph{superbasis}. A subbasis is defined to be a
truncated orthonormal basis, tending towards completeness, which has been
variationally optimized to minimize the energy of one particular state.
Therefore, we expect a loose association between a subbasis and an energy
eigenstate. The superbasis is a formed from the union of multiple subbases each
independently optimized to different eigenstates. The superbasis is no longer
orthonormal, as each subbasis is allowed to vary its variational parameters
independently. As well, the superbasis tends towards overcompleteness as more
states are added. For a numerical method such as ours, we do need to be
concerned with numerical instability or intractability introduced by linear
dependence and overcompleteness within the superbasis. Despite this, we find
that the variational procedure applied to each subbasis helps to ameliorate
these problems with numerical instability. As a potential explanation, recall
that each subbasis can be loosely associated with each energy eigenstate, and
these eigenstates are orthogonal to each other. Serendipitously then, it seems
the process of variational optimization mutually orthogonalizes the subbases and
reduces numerical instability.